% Generated draft: requires mathematical and source review.
\documentclass[11pt]{article}
\usepackage[margin=25mm]{geometry}
\usepackage{fontspec}
\setmainfont{Latin Modern Roman}
\usepackage{amsmath,amssymb,mathtools,unicode-math}
\setmathfont{Latin Modern Math}
\usepackage{enumitem,needspace,xurl,hyperref}
\setlength{\parindent}{0pt}
\setlength{\parskip}{5pt}
\setlength{\emergencystretch}{3em}
\newcommand{\N}{\mathbb{N}}
\newcommand{\Z}{\mathbb{Z}}
\newcommand{\Q}{\mathbb{Q}}
\newcommand{\R}{\mathbb{R}}
\newcommand{\C}{\mathbb{C}}
\newcommand{\F}{\mathbb{F}}
\newcommand{\A}{\mathbb{A}}
\newcommand{\PP}{\mathbb{P}}
\newcommand{\E}{\mathbb{E}}
\newcommand{\eps}{\varepsilon}
\newcommand{\dd}{\,\mathrm{d}}
\begin{document}
\section*{UnsolvedMath ID: 30004273}
\textbf{Source problem code:} catalogue entry
\subsection{Definitions and mathematical statement}
For the concrete model, let
\[
f(x,z)=z-x^2,\qquad \Sigma=\{(x,z)\in\R^2:z=x^2\},\qquad M=\R^2\setminus\Sigma.
\]
On \(M\), equip the tangent bundle with the Riemannian metric for which
\[
X_1=\partial_x,\qquad X_2=f\partial_z
\]
are orthonormal. Thus
\[
g=\dd x^2+\frac{\dd z^2}{(z-x^2)^2},\qquad \dd\mu_g=\frac{\dd x\,\dd z}{|z-x^2|}.
\]
The singular set is where the two vector fields become linearly dependent. Along \(\Sigma\), their span is \(\operatorname{span}\{\partial_x\}\); at \((0,0)\), this span equals \(T_{(0,0)}\Sigma\), so the origin is a tangency point.

To make the operator question precise, use the standard minimal realization on the complex Hilbert space \(H=L^2(M,\mu_g;\C)\):
\[
\mathcal D(A_0)=C_c^\infty(M;\C),\qquad A_0u=-\Delta_g u,
\]
where
\[
\Delta_g u
=|f|\,\partial_x\!\left(|f|^{-1}\partial_xu\right)
 +|f|\,\partial_z\!\left(|f|\partial_zu\right).
\]
Here \(C_c^\infty(M;\C)\) denotes smooth functions with compact support contained in the regular region. This initial domain imposes no boundary or transmission condition on \(\Sigma\). Essential self-adjointness means that the closure \(\overline{A_0}\) is self-adjoint in \(H\), equivalently \(\overline{A_0}=A_0^*\).

\textbf{Concrete question.} Is \(A_0\) essentially self-adjoint on all of \(M=\R^2\setminus\Sigma\)? The negative sign in the definition of \(A_0\) does not affect this question.

\textbf{Broader source question and unresolved qualification.} The source also asks whether the Laplace--Beltrami operator on the regular region of an almost-Riemannian structure is essentially self-adjoint when its singular set has tangency points. It does not specify a quantified class of structures or the global hypotheses for that broader question, so a precise universal assertion cannot be extracted. The Hilbert space, orthonormal-frame interpretation, and minimal initial domain used above are standard interpretations of the concrete question, but are not explicitly stated in the supplied record and require confirmation against the source.
\subsection{Short English statement}
For the metric making \(\partial_x\) and \((z-x^2)\partial_z\) orthonormal away from the parabola \(z=x^2\), ask whether the Laplace--Beltrami operator initially defined on compactly supported smooth functions has a unique self-adjoint extension in the Riemannian \(L^2\) space. The origin is a tangency point. The broader question about arbitrary almost-Riemannian structures lacks specified global hypotheses.
\subsection{Sources}
\begin{itemize}
\item Supplied local UnsolvedMath registry record, permanent ID \(30004273\), title \emph{Self-Adjointness Across Almost-Riemannian Tangency Points}, source problem code \texttt{catalogue entry}, in the Oberwolfach Reports Open Problems collection. The record lists proposed year \(2019\) and supplies no individual proposer.
\item Supplied source citation: \emph{Mini-Workshop: Self-adjoint Extensions in New Settings (2020)}, Oberwolfach Reports, pp.~2911--2950, DOI \texttt{10.4171/owr/2019/47}. \url{https://doi.org/10.4171/owr/2019/47}. The background describes Oberwolfach Reports as documenting mathematical workshops and their open-problem discussions.
\item Supplied literature source, labeled ``OWR 2019 47'': \url{https://ems.press/content/serial-article-files/46828}.
\item Supplied literature source, labeled ``Beschastnyi'': \url{https://doi.org/10.1007/s00025-022-01832-z}. No article title or fuller author attribution is supplied in the record.
\item Literature assessment checked at \texttt{2026-08-22}: the supplied curation labels the problem open, identifies the tangency-point parabola model as the hardest unresolved case, and states that results on operator closure and domains do not establish essential self-adjointness for this model. This is supplied dated curation derived from the OWR material, not an independently verified assessment. The record expressly calls for independent verification of time-sensitive status and cited claims.
\end{itemize}
\end{document}
